\section{Conclusiones}

El análisis de PuellaGame permitió transformar una necesidad operativa urgente
en un conjunto de requerimientos verificables. La organización necesita
conservar información de sucursales, premios y clientes, recuperarla mediante
una llave y modificarla sin perder persistencia. El prototipo responde a ese
alcance mediante menús de consola, operaciones CRUD, validación de entradas,
manejo de excepciones y un archivo CSV independiente para cada entidad.

La separación entre modelos, repositorios, servicios, interfaz y validaciones
resultó fundamental para que el almacenamiento en archivos no se confundiera
con la lógica de negocio. Esta arquitectura permite probar cada responsabilidad
de forma independiente y deja localizada la sustitución futura de los
repositorios CSV por componentes conectados a una base de datos.

La comparación tecnológica muestra que los CSV son apropiados como mecanismo
temporal: son simples, legibles y suficientes para iniciar la captura. No
obstante, no proporcionan por sí mismos integridad centralizada, consultas
declarativas, concurrencia, recuperación ni administración detallada de
usuarios. Cuando PuellaGame evolucione hacia una operación con mayor volumen y
múltiples usuarios, una base de datos será la alternativa adecuada. En este
sentido, el prototipo no sustituye el diseño posterior, sino que produce un
primer acervo estructurado y validado para facilitarlo.
